\documentclass{article}

% if you need to pass options to natbib, use, e.g.:
    \PassOptionsToPackage{numbers, compress}{natbib}
% before loading neurips_2024


% camera-ready / non-anonymized
\usepackage[final]{neurips_2024}


% to compile a preprint version, e.g., for submission to arXiv, add add the
% [preprint] option:
%     \usepackage[preprint]{neurips_2024}


% to compile a camera-ready version, add the [final] option, e.g.:
%     \usepackage[final]{neurips_2024}


% to avoid loading the natbib package, add option nonatbib (load neurips_2024 only once)
%   \usepackage[final,nonatbib]{neurips_2024}


\usepackage[utf8]{inputenc} % allow utf-8 input
\usepackage[T1]{fontenc}    % use 8-bit T1 fonts
\usepackage{hyperref}       % hyperlinks
\usepackage{url}            % simple URL typesetting
\usepackage{booktabs}       % professional-quality tables
\usepackage{graphicx}       % for resizing tables/figures
\usepackage{amsmath}        % aligned equations (align, gather, ...)
\usepackage{amssymb}        % extra math symbols (e.g., \varnothing)
\usepackage{amsfonts}       % blackboard math symbols
\usepackage{nicefrac}       % compact symbols for 1/2, etc.
\usepackage{microtype}      % microtypography
\usepackage{xcolor}         % colors


\title{When Senses Disagree: Calibrated Arbitration for Conflicting Multimodal Reasoning}


% The \author macro works with any number of authors. There are two commands
% used to separate the names and addresses of multiple authors: \And and \AND.
%
% Using \And between authors leaves it to LaTeX to determine where to break the
% lines. Using \AND forces a line break at that point. So, if LaTeX puts 3 of 4
% authors names on the first line, and the last on the second line, try using
% \AND instead of \And before the third author name.


\author{%
  Otto Xin \\
  Department of Communication Studies\\
  Northwestern University\\
  Evanston, IL 60208 \\
  \texttt{haohangxin@u.northwestern.edu} \\
  \And
  Nicholas Ornstein \\
  Department of Communication Studies\\
  Northwestern University\\
  Evanston, IL 60208 \\
  \texttt{nickornstein@u.northwestern.edu}
}



\begin{document}


\maketitle

\begin{abstract} % Renamed the abstract heading
  Multimodal large language models (MLLMs) have achieved strong performance on visual question answering and instruction following, yet they remain brittle when visual and textual inputs conflict. In these modality-conflict settings, models can generate fluent but ungrounded outputs and exhibit unstable evidence use, inconsistently privileging visual cues or textual priors. We address this reliability gap by introducing \textbf{CARM} (Conflict-Aware Reasoning Module), a lightweight, model-agnostic arbitration layer that makes conflict handling explicit. Given an image, a caption/instruction, and a question, CARM predicts a typed conflict signal, estimates question-conditioned modality reliabilities via unimodal probes, and selects an action that governs response generation (trust vision, trust text, require agreement, or abstain). The goal is to reduce incorrect yet confident answers under modality contradiction while preserving performance on consistent inputs, and to expose interpretable intermediate signals for diagnosing model behavior.

  We evaluate CARM on a controlled Conflict Suite derived from clean image--text pairs with held-out splits to assess generalization. Evaluation goes beyond end-task accuracy to include conflict detection, arbitration quality, calibration, and selective prediction under abstention, with comparisons to standard backbone and prompting-based baselines. Deliverables include the CARM implementation, the dataset and split specification, a unified training and evaluation pipeline, and a reproducibility package for end-to-end reruns.

\end{abstract}

\section{Goal}
\label{sec:goal}

This project aims to improve the reliability of multimodal large language models (MLLMs) when visual and textual inputs provide conflicting information. We propose a lightweight, model-agnostic "Conflict-Aware Reasoning Module" (CARM), that makes cross-modal arbitration explicit by classifying conflict type, estimating question-conditioned modality reliability, and selecting an action that governs how the model responds (e.g., trust vision, trust text, require agreement, or abstain). The overarching goal is to reduce incorrect yet fluent answers under contradiction while preserving performance on non-conflicting inputs, and to provide interpretable intermediate signals that enable principled evaluation of robustness, calibration, and selective prediction behavior.

\section{Background and Motivation}
\label{sec:background}

Recent multimodal large language models (MLLMs) have enabled general-purpose visual question answering and instruction following by coupling a pretrained vision encoder with a large language model through a lightweight cross-modal interface, often while freezing most backbone parameters to reduce training cost \citep{dai_instructblip_2023, li_blip-2_2023, liu_visual_2023}. This design has proven highly effective across a wide range of tasks, but it also introduces well-documented reliability challenges. In particular, MLLMs have been shown to confidently return text that conflicts with the visual input, a phenomenon widely studied under the umbrella of hallucination and faithfulness errors \citep{li_evaluating_2023, favero_multi-modal_2024, seth_towards_2025}. Such failures are not merely cosmetic as they can undermine downstream decision-making and safety-critical use cases where users implicitly assume that model outputs are grounded in the provided evidence.

A principal source of these errors is that real-world multimodal inputs are frequently internally inconsistent. This may be due factors including the divergence of images, captions, and user-provided context, stale or incorrect metadata, adversarial manipulation, or missing visual evidence. This phenomenon has been formalized as modality conflict, in which a model must arbitrate between contradictory cues across modalities \citep{zhang_robust_2025}. Empirical evidence indicates that such conflicts induce substantial performance degradation even when each modality is, in isolation, informative, implying that failures often arise from deficient cross-modal fusion and arbitration rather than perception alone \citep{zhang_robust_2025, hua_how_2025}. Large-scale evaluations further suggest that contemporary MLLMs lack an explicit conflict-resolution mechanism in that under contradictory image--caption pairs, leading models (e.g., LLaVA, InstructBLIP, Qwen-VL) exhibit unstable evidence use, alternately privileging visual or textual cues in a model- and instance-dependent manner \citep{hua_how_2025}. Motivated by these findings, recent benchmarks have begun to directly operationalize conflict sensitivity. For example, CLASH evaluates cross-modal contradiction detection using controlled object- and attribute-level contradictions and reports that many open models struggle to reliably identify inconsistencies \citep{popordanoska_clash_2025}. Complementary work extends this line of evaluation to additional modalities and contexts, including layout-rich inconsistency reasoning and broader hallucination taxonomies \citep{yan_multimodal_2025, chen_survey_2025}.

Beyond detecting conflict, recent studies emphasize that MLLMs exhibit systematic modality preferences. When modalities disagree, models often default to one input source, and the direction and strength of this bias varies by model family and context \citep{zhang_evaluating_2025, hua_how_2025}. This modality dominance is not purely a property of the task; it can be partially shifted through prompting and representation-level interventions, indicating that arbitration behavior is, at least in part, controllable \citep{zhang_evaluating_2025}. At the same time, apparent “preference” is intertwined with unimodal competence. Models may seem to follow one modality because they are more uncertain when reasoning from the other in general. Recent work formalizes this distinction by decomposing modality-bias into relative unimodal reasoning uncertainty and an inherent preference term, and reports monotonic relationships between uncertainty gaps and the probability of following a given modality under conflict \citep{zhang_when_2025}. Evidence from human-AI cross-modal perception comparisons further highlights that such biases can diverge from human integration strategies. In settings where human participants reliably prioritize auditory cues over misleading visual signals, AI systems can remain strongly vision-dominant and fail to suppress irrelevant visual information \citep{jia_seeing_2025}. Mechanistic analyses suggest that these arbitration behaviors are implemented through specific internal pathways; in particular, instruction tokens appear to function as structural anchors through which multimodal cues are routed, with competition resolved in deeper layers and a small set of attention heads exerting disproportionate causal influence \citep{zhang_instruction_2026}. These findings motivate conflict-handling methods that are explicitly uncertainty-aware and aligned with where arbitration occurs in the network, rather than treating conflict as an undifferentiated end-task fine-tuning problem.

Existing mitigation strategies have largely fallen into three categories: prompting-based verification, supervised fine-tuning for robustness, and reinforcement learning style optimization, often evaluated on conflict-oriented benchmarks \citep{zhang_evaluating_2025}. While these approaches can reduce hallucination rates, they may be difficult to interpret, they can be sensitive to the training distribution, and often conflate multiple behaviors (conflict detection, modality choice, and answer generation) into a single black-box improvement. In parallel, abstention has been proposed as a general safety mechanism for language models, providing a controllable way to trade off coverage for reduced error under uncertainty \citep{wen_know_2025}. However, abstention in multimodal settings remains under-specified. A model should abstain not merely when it is uncertain, but when available evidence is inconsistent or too unclear for the specific question.

We aim to make cross-modal arbitration explicit and evaluable by introducing CARM, a lightweight module that (1) detects and types conflicts, (2) estimates question-conditioned modality reliability using unimodal probes and counterfactual modality corruption, and (3) selects an action policy that governs evidence usage and permits calibrated abstention. This design is intended to improve robustness under contradiction while preserving performance on consistent inputs, and to yield intermediate signals that support diagnostic evaluation beyond end-task accuracy \citep{popordanoska_clash_2025, zhang_evaluating_2025, zhang_instruction_2026, wen_know_2025}.


\section{Architecture}
\label{sec:architecture}

\subsection{Overview}
We propose \textbf{CARM} (Conflict-Aware Reasoning Module), a lightweight overlay that turns implicit modality preference into an explicit, testable arbitration process. Given image $I$, text $T$ (caption or instruction), and question $q$, CARM outputs intermediate signals for (1) conflict type, (2) question-conditioned modality reliability, and (3) a discrete action that determines how the final answer is generated.

\subsection{Backbone model and interface}
CARM wraps an existing open MLLM as a frozen backbone. The backbone produces token-level hidden states
\[
H = \mathrm{Backbone}(I, T, q), \qquad H \in \mathbb{R}^{L \times d}.
\]
We define a subset of \emph{anchor tokens} $A$ corresponding to the question and instruction spans, and denote their states by $H_A$. CARM reads from these anchor states, since arbitration should be conditioned on representations most directly tied to the question being answered.

\subsection{Unimodal probes for evidence and uncertainty}
To make arbitration question-dependent (and reduce collapse to a static modality prior), we compute two controlled unimodal probe passes:
\begin{align}
p_v(y \mid I, q) &= \mathrm{Backbone}(I, \varnothing, q), \\
p_t(y \mid T, q) &= \mathrm{Backbone}(\varnothing, T, q),
\end{align}
where $\varnothing$ denotes the suppressed modality (for example, omitting the caption or masking visual tokens).

From each probe distribution we extract a compact feature vector:
\[
\phi_v = \Phi(p_v), \qquad \phi_t = \Phi(p_t),
\]
including simple uncertainty proxies such as entropy, margin between top candidates, and (optionally) variance under repeated sampling. These probe features summarize whether the backbone can answer confidently from each modality in isolation.

\subsection{CARM heads: conflict, reliability, and action}
CARM contains three lightweight prediction heads. We first pool anchor states into a fixed-dimensional vector
\[
h = \mathrm{Pool}(H_A),
\]
where $\mathrm{Pool}(\cdot)$ can be mean pooling or attention pooling over anchor tokens. We then form a joint decision input
\[
u = [h;\ \phi_v;\ \phi_t],
\]
and predict (i) conflict type, (ii) modality reliabilities, and (iii) an action policy:
\begin{align}
  \hat{c} &= \mathrm{softmax}(W_c u), \\
  \hat{r}_v, \hat{r}_t &= \sigma(W_r u), \\
  \hat{a} &= \mathrm{softmax}(W_a u).
\end{align}

We use a typed conflict label space
\[
c \in \{\texttt{none},\ \texttt{object},\ \texttt{attribute},\ \texttt{relation},\ \texttt{count}\},
\]
and a discrete action space
\[
a \in \{\texttt{trust-vision},\ \texttt{trust-text},\ \texttt{require-agreement},\ \texttt{abstain}\}.
\]
This separation allows evaluation of arbitration quality independently from answer accuracy.

\subsection{Action-conditioned generation}
The action $a$ controls the evidence available at generation time:
\begin{itemize}
  \item \texttt{trust-vision}: generate using $(I,q)$ while suppressing $T$.
  \item \texttt{trust-text}: generate using $(T,q)$ while suppressing $I$.
  \item \texttt{require-agreement}: if unimodal probe predictions agree (under a task-specific equivalence rule), generate the shared answer; otherwise abstain.
  \item \texttt{abstain}: emit a standardized refusal (or a short clarifying request, if allowed).
\end{itemize}
Suppression can be implemented by omitting the corresponding modality from the input or masking its tokens. The key requirement is that the final response respects the arbitration decision.

\subsection{Training objective}
We freeze the backbone and train only the CARM heads with a multi-task objective:
\[
\mathcal{L} = \mathcal{L}_{\text{type}} + \mathcal{L}_{\text{act}} + \mathcal{L}_{\text{rel}} + \lambda\, \mathcal{L}_{\text{cf}}.
\]
$\mathcal{L}_{\text{type}}$ and $\mathcal{L}_{\text{act}}$ are cross-entropy losses for conflict type and action prediction. $\mathcal{L}_{\text{rel}}$ supervises $\hat{r}_v,\hat{r}_t$ (for example, via regression or binary cross-entropy against reliability targets).

To prevent the reliability estimator from collapsing to a dataset-level modality prior, we introduce a counterfactual constraint $\mathcal{L}_{\text{cf}}$ built from paired examples where exactly one modality is corrupted. For a text-corrupted variant $T^{\dagger}$ (caption swap, negation injection, or attribute flip) with $I$ fixed, we encourage text reliability to drop:
\[
\mathcal{L}_{\text{cf}}^{(t)} = \max\!\left(0,\ m - \left(\hat{r}_t(I,T,q) - \hat{r}_t(I,T^{\dagger},q)\right)\right).
\]
For a vision-corrupted variant $I^{\dagger}$ (blur, occlusion, or distractor replacement) with $T$ fixed, we analogously encourage vision reliability to drop:
\[
\mathcal{L}_{\text{cf}}^{(v)} = \max\!\left(0,\ m - \left(\hat{r}_v(I,T,q) - \hat{r}_v(I^{\dagger},T,q)\right)\right).
\]
We set $\mathcal{L}_{\text{cf}} = \mathcal{L}_{\text{cf}}^{(t)} + \mathcal{L}_{\text{cf}}^{(v)}$.


\section{Experiments}
\label{sec:experiments}

\subsection{Evaluation questions}
The experimental design targets four evaluation questions. First, does CARM improve answer quality in the presence of image--text conflict while preserving performance on consistent inputs? Second, does CARM select appropriate arbitration actions (trusting vision, trusting text, requiring agreement, or abstaining) in a way that aligns with the underlying evidence required by the question? Third, are the predicted modality reliabilities calibrated and responsive to targeted degradations of the corresponding modality? Fourth, do observed gains generalize beyond the conflict patterns encountered during training, rather than reflecting overfitting to a particular construction procedure?

\subsection{Data construction and splits}
\paragraph{Conflict Suite.}
We construct a controlled Conflict Suite from clean image--text pairs and associated questions. Starting from consistent instances, we generate conflict cases via (1) caption swaps, including semantically similar (same-topic) swaps to reduce trivial topic mismatch cues; (2) text edits that induce typed contradictions, focusing on attribute-level changes (e.g., color, count) and relational changes (e.g., left/right); and (3) vision corruptions (blur, occlusion, and distractor replacement) at multiple severity levels. In addition, we include hard non-conflict examples that are semantically close yet consistent, to discourage shortcut detection based on superficial distributional cues.

\paragraph{Held-out generalization.}
To support generalization claims, we define evaluation splits prior to training. We reserve at least one entire conflict family (e.g., relation edits) as a held-out test set that is never observed during training. Separately, for vision corruptions, we evaluate on corruption severities stronger than those used for training. When feasible, we further evaluate transfer on one to two external benchmarks that probe contradiction detection and/or reliability under conflict, using metrics compatible with each benchmark’s annotation format.

\subsection{Baselines}
We compare CARM against baselines that isolate the role of training, prompting, and abstention. The primary baselines include: (i) the frozen backbone with standard decoding on $(I,T,q)$; (ii) a prompt-based verification strategy that explicitly requests consistency checking prior to answering; (iii) an uncertainty-threshold abstention baseline that abstains when the backbone’s uncertainty exceeds a tuned threshold; and (iv) a probe-only heuristic that selects between vision-only and text-only answering based solely on unimodal probe uncertainty features, without a learned arbitration head.

\subsection{Metrics}
We report end-task performance separately on consistent and conflict subsets, with additional breakdowns by conflict family. For intermediate behavior, we evaluate conflict-type prediction using macro-F1 and per-type F1. We evaluate arbitration via \emph{action accuracy} over $a \in \{\texttt{trust-vision},\ \texttt{trust-text},\ \texttt{require-agreement},\ \texttt{abstain}\}$, where the reference action is determined by construction (which modality is corrupted) together with question type (which modality contains the required evidence). To evaluate abstention as a first-class outcome, we report risk--coverage curves that characterize the accuracy--coverage trade-off as abstention thresholds vary. Finally, we evaluate reliability calibration using binned calibration plots and targeted corruption analyses: as corruption severity increases for a given modality, the corresponding reliability should decrease monotonically.

\subsection{Ablations and robustness checks}
We conduct ablations to test whether performance improvements arise from the intended design components. Specifically, we compare heads attached to question/instruction anchor states versus pooled representations, remove the counterfactual constraint $\mathcal{L}_{\mathrm{cf}}$, remove unimodal probe features $(\phi_v,\phi_t)$, and remove abstention by restricting the action space to always answer. Robustness checks include same-topic caption swaps and paraphrase-invariant variants of $T$, which help distinguish genuine contradiction reasoning from superficial cue exploitation. We also evaluate beyond the corruption severities seen during training to test whether reliability behaves consistently under increasing modality degradation.

\section{Deliverables}
\label{sec:deliverables}

We will produce the following deliverables to support empirical evaluation and reproducibility.

\begin{enumerate}
  \item \textbf{CARM implementation.} A modular implementation of CARM that wraps a frozen open-source MLLM backbone, including (1) unimodal probe passes, (2) conflict-type, reliability, and action heads, and (3) action-conditioned generation with standardized abstention behavior.

  \item \textbf{Conflict Suite dataset and split specification.} A curated dataset of clean and conflict instances with typed conflict labels, corruption metadata (modality and severity), and pre-defined held-out generalization splits (held-out conflict family and held-out corruption severity). We will provide documentation of the construction pipeline and split criteria.

  \item \textbf{Unified training and evaluation pipeline.} Scripts that train CARM heads on fixed backbones and evaluate all methods (CARM and baselines) under identical protocols, producing per-example logs of intermediate outputs (conflict type, $r_v$, $r_t$, action, abstention) and final answers.

  \item \textbf{Benchmarking outputs.} Tables and figures summarizing (1) accuracy on consistent versus conflict subsets with per-family breakdowns, (2) conflict detection performance, (3) action accuracy, (4) calibration analyses (including corruption-strength curves)

  \item \textbf{Reproducibility package.} Environment specifications, fixed random seeds, and configuration files sufficient to reproduce the main results and plots, along with a short guide to rerun experiments end-to-end.
\end{enumerate}



\section{Team Organization}
\label{sec:team}

The project will be conducted by Otto and Nick with shared responsibility for core design decisions, implementation, and analysis. Otto will lead the methodological development of CARM, including the architecture implementation (unimodal probes, prediction heads, and action-conditioned generation) and training pipeline integration with the selected backbone. Nick will lead dataset construction and experimental evaluation, including the Conflict Suite generation, held-out split design, baseline implementation, and metric computation (conflict-type performance, action accuracy, calibration analyses, and risk--coverage curves). Both members will jointly conduct ablation studies, error analysis, and writing, with coordinated weekly integration checkpoints to ensure compatibility between the data pipeline and model evaluation framework.


\section{Time Table}
\label{sec:timetable}

\paragraph{Week 6 (Design and setup).}
Finalize the task formulation and experimental protocol, including the Conflict Suite construction rules and the held-out generalization splits. Select the backbone model and implement the full CARM pipeline skeleton (unimodal probes, heads, action-conditioned generation), together with unified logging and evaluation utilities.

\paragraph{Week 7 (Data construction and baselines).}
Generate the initial Conflict Suite (clean + conflict families + corruption severities) and verify data quality. Implement and run core baselines (backbone-only, prompt-based verification, uncertainty-threshold abstention, probe-only heuristic) to establish reference performance and identify dominant failure modes.

\paragraph{Week 8 (Training and primary evaluation).}
Train CARM heads on the frozen backbone using the multi-task objective. Run primary evaluations on in-distribution and held-out generalization splits, and produce the main result tables (accuracy by subset/family; conflict-type F1; action accuracy).

\paragraph{Week 9 (Ablations, calibration, and write-up).}
Complete required ablations (attachment site, removing probes, removing counterfactual constraint, removing abstention). Produce calibration analyses (binned calibration plots; corruption-strength curves) and risk--coverage curves. Consolidate results and finalize the proposal manuscript sections and figures.

\paragraph{Week 10 (Presentation).}
Prepare presentation slides, including the architecture figure, experimental setup overview, main quantitative results, and a concise error analysis illustrating representative conflict cases and arbitration behavior.


\begin{table}[t]
\centering
\caption{Baseline comparison on \texttt{test\_id} (native behavior). We report Coverage (fraction answered), Accuracy on answered examples, and Task Success (policy-aware success). Category-wise Task Success highlights modality-routing cases (C3: text irrelevant; C4: vision irrelevant). Direct never abstains, so it is expected to score poorly on abstention-target categories (C2/C5).}
\label{tab:baseline-main}
\setlength{\tabcolsep}{3pt}
\scriptsize
\resizebox{\columnwidth}{!}{%
\begin{tabular}{lcccccc}
\toprule
Method & Coverage $\uparrow$ & Acc. (answered) $\uparrow$ & TaskSuccess (overall) $\uparrow$ & TaskSuccess C1 $\uparrow$ & TaskSuccess C3 $\uparrow$ & TaskSuccess C4 $\uparrow$ \\
\midrule
Direct (\texttt{backbone\_direct}) & 1.00 & 0.78 & 0.52 & 0.80 & 0.76 & 0.74 \\
Verification (\texttt{prompt\_verification}) & 0.62 & 0.84 & 0.61 & 0.87 & 0.41 & 0.39 \\
Unc.-Thresh Abstain (\texttt{uncertainty\_threshold\_abstain}) & 0.80 & 0.82 & 0.66 & 0.83 & 0.55 & 0.52 \\
Probe Heuristic (\texttt{probe\_only\_heuristic}) & 0.78 & 0.85 & 0.72 & 0.86 & 0.68 & 0.66 \\
\bottomrule
\end{tabular}
}
\end{table}


%\section*{References}


% References follow the acknowledgments in the camera-ready paper. Use unnumbered first-level heading for
% the references. Any choice of citation style is acceptable as long as you are
% consistent. It is permissible to reduce the font size to \verb+small+ (9 point)
% when listing the references.
% Note that the Reference section does not count towards the page limit.
% \medskip



%%%%%%%%%%%%%%%%%%%%%%%%%%%%%%%%%%%%%%%%%%%%%%%%%%%%%%%%%%%%

%\appendix

%\section{Appendix / supplemental material}

%%%%%%%%%%%%%%%%%%%%%%%%%%%%%%%%%%%%%%%%%%%%%%%%%%%%%%%%%%%%

\newpage


\bibliographystyle{plainnat}
\bibliography{cs396_zotero,ref}

\end{document}
